\section{An approximate transfer theorem for formal-arc orders}
\label{V:sec-transfer}

An order along a generically dominant formal arc records vanishing of
effective Cartier divisors even when the arc's center is not a closed
point.  We will transfer bounded orders in a twisted sequence to
arbitrarily small normalized orders in anticanonical systems.

\begin{theorem}[Approximate arc transfer]\label{V:approximate-arc}
Let $W$ be smooth connected projective over $\mathbb C$, with
$q(W)=0$, and let $A=-K_W$ admit a smooth semipositive Hermitian
metric.  Let $M$ be a line bundle with $-M$ pseudoeffective.  Let $v$
be the order along a formal arc over an extension field of
$\mathbb C$ whose fraction point maps to the generic point of $W$;
the trivial order is allowed.  Suppose there are effective integral
divisors
\[
 N_i\sim M+m_iA,\qquad 0<m_1<m_2<\cdots,
 \qquad v(N_i)=O(1).
\]
For every real $\delta>0$ there are an integer $r>0$ and an effective
integral divisor $N\sim rA$ with $v(N)<\delta r$.
\end{theorem}

The exact scalar theorem proved earlier yields a stronger conclusion
under these hypotheses.  The proof here explains a different
mechanism.  A fixed part absorbs none of the prescribed arc order;
a nef reduction then leaves a base line whose first jets can be
extended.  Bigness of that line supplies sections with arbitrarily
small normalized order.  Both the assumption $q(W)=0$ and the
pseudoeffectivity of $-M$ have specific roles in this construction.

\begin{proof}[Proof of Theorem~\ref{V:approximate-arc}]
The assertion is immediate if $W$ is a point.  We assume positive
dimension.  The arc lifts uniquely to every proper birational model,
by the valuative criterion.  We use the same letter $v$ for orders
there; orders of effective rational Cartier divisors are finite and
nonnegative, and are preserved by pullback.

\subsection*{The common field and its rank-one fibers}

Call the indices $i\ge2$ interior.  For each $i$, let $K_i$ be the
subfield of $\mathbb C(W)$ generated by ratios of nonzero sections
of the same positive multiple of $N_i$.  We first show that $K_i$
is independent of the interior index.

For an interior $i$ and any $j$, there is a rational $c>0$ such that
$N_i-cN_j$ is rationally linearly equivalent to an effective rational
divisor.  For $j<i$ choose an index beyond $i$, and for $j>i$ choose
an earlier index, then express $m_i$ as the corresponding strict
rational convex combination.  The same combination of the effective
divisors has class $M+m_iA$.  The case $i=j$ is immediate.  Given
a ratio of sections of $kN_j$, choose a sufficiently divisible $d$
so that $dN_i-kN_j$ has a nonzero section.  Multiplying numerator
and denominator by that section realizes the same ratio in $|dN_i|$.
Thus $K_j\subset K_i$, and the fields for interior indices coincide;
write $K$ for their common value.

Choose finitely many of these ratios attaining the transcendence
degree of $K$ and place them in one complete system $|dN_{i_0}|$.
Resolve its map and take the Stein factor.  Its function field $F$
is relatively algebraically closed in $\mathbb C(W)$ and contains
$K$: the field of the selected image is contained in $K$ with the
same transcendence degree, so every element of $K$ is algebraic over
it.  Apply Lemma~\ref{V:flat-model} and then resolve to obtain
\[
 S\longrightarrow V,\qquad f:V\longrightarrow Y,
 \qquad p:V\longrightarrow W,
\]
where $Y,S$ are smooth projective, $V$ is normal equidimensional over
$Y$, $p$ is birational, and $\mathbb C(Y)=F$.  We also write $f,p$
for the maps from $S$.  Set $A_S=p^*A$ and
$E=K_S-p^*K_W\ge0$.  The divisor $E$ is exceptional.  Thus, for a
line bundle $Q$ on $W$ and an integer $u\ge0$,
\[
 H^0(S,p^*Q+uE)=s_E^u p^*H^0(W,Q),
\]
by normality of $W$ and extension across codimension two.

The pullback of the hyperplane bundle from the selected image is big
on its Stein factor and remains big after modifying the base.  It is
linearly dominated by $d p^*N_{i_0}$.  The interpolation above therefore
shows that a sufficiently divisible multiple of any interior
$p^*N_i$ linearly dominates the pullback of any fixed divisor on $Y$.
This remains true on further base modifications.

For every $l>0$ and every interior $i$, on the generic fiber one has
\begin{equation}\label{V:generic-rank-one}
 h^0(S_\eta,l(p^*N_i+E)|_{S_\eta})
 =h^0(S_\eta,lp^*N_i|_{S_\eta})=1.
\end{equation}
Both spaces contain the indicated effective section.  If either had
two independent sections, a sufficiently ample base twist would extend
them globally by generation of the twisted direct image.  The preceding
domination clears this twist using a common effective multiplier.
Their ratio would then be a ratio in a multiple of $p^*N_i+E$ and
would not belong to $F$, by independence over the generic base field.
Exceptional pushforward identifies its ratio field with that of $N_i$,
contradicting $K_i\subset F$.

Rank one now turns interpolation of line bundles into interpolation
of horizontal divisors.  If interior indices $i<j<k$ satisfy
$m_j=c m_i+(1-c)m_k$ with $0<c<1$ rational, then
\[
 N_j\sim_{\mathbb Q}cN_i+(1-c)N_k.
\]
Clear denominators and restrict to the generic fiber.  The two effective
sections belong to the same one-dimensional space in
\eqref{V:generic-rank-one}; their ratio is a scalar in $F$ and has
no horizontal divisor.  Consequently
\[
 (p^*N_j)^{\mathrm{hor}}
 =c(p^*N_i)^{\mathrm{hor}}+(1-c)(p^*N_k)^{\mathrm{hor}}.
\]
Fixing two interior indices gives rational horizontal divisors $D_0,C_0$
with fixed finite support such that
\[
 (p^*N_i)^{\mathrm{hor}}=m_iD_0+C_0\quad(i\ge2).
\]
Since the left side is effective for arbitrarily large $m_i$, each
coefficient of $D_0$ is nonnegative.

Take two interior indices $j>i$ and set $G=p^*(N_j-N_i)$ on $V$.
Its horizontal part is nonnegative.  For each prime divisor $Q$ of $Y$
take the minimum of
\[
 \frac{\operatorname{coeff}_{\Gamma}G}
      {\operatorname{coeff}_{\Gamma}f^*Q}
 \quad\hbox{over the components $\Gamma$ dominating $Q$.}
\]
There are only finitely many nonzero minima.  Subtracting their
pullback makes $G$ effective: the equidimensional normal model has
no other vertical primes.  After clearing denominators this gives
\begin{equation}\label{V:base-decomposition}
 a p^*A\sim D^V+f^*B,\qquad a>0,
\end{equation}
where $D^V$ is effective Cartier, $B$ is Cartier, and over every prime
of $Y$ some component has coefficient zero in $D^V$.  On $S$ put
$D=(D^V)|_S$ by pullback.  Its horizontal part is $aD_0$.

For fixed $l\ge0$ and sufficiently large $i$, the horizontal formula
gives an effective generic-fiber divisor
\[
 R_{i,l}=(p^*N_i-lD)|_{S_\eta}\ge0.
\]
Since $aA_S|_{S_\eta}\sim D|_{S_\eta}$, multiplication by
$s_{R_{i,l}}$ embeds the section spaces of $laA_S$ and $laA_S+E$
on the generic fiber into those of $p^*N_i$ and $p^*N_i+E$,
respectively.  The latter spaces have dimension one by
\eqref{V:generic-rank-one}.  The former contain $s_D^l$ and
$s_Es_D^l$ in base frames, so they too have dimension one.  For $l=0$
this also uses connectedness of the generic fiber.
In particular,
\begin{equation}\label{V:global-sections-base}
 H^0(W,laA)=s_D^l f^*H^0(Y,lB)
\end{equation}
under pullback to $S$.  Indeed any section has this form with a
rational base section by the generic-fiber statement.  On $V$, a pole
over a base prime would remain a pole on its coefficient-zero
component, which is impossible.  The rational base section is therefore
regular.  Conversely every section on the right is regular on $S$ and
descends to $W$.

\subsection*{Removing the fixed divisor from the arc order}

Pass to an infinite subsequence with $m_i$ in one residue class modulo
$a$, and fix an index $j$ in it.  For subsequent indices put
$l_i=(m_i-m_j)/a$.  Horizontal equality and
\eqref{V:base-decomposition} give actual divisor identities on $V$
\[
 p^*N_i=p^*N_j+l_iD^V+f^*V_i,
 \qquad V_i\sim l_iB.
\]
To justify the base divisor, subtract first $l_iB$ using the line
bundle identification.  The remaining principal divisor is vertical;
its rational function has zero divisor on the proper normal generic
fiber, and hence belongs to its constant field $F$.

The negative parts of the $V_i$ are bounded by a fixed effective
Cartier divisor $R$ on $Y$.  For a negative coefficient over $Q$, use
a component where $D^V$ has coefficient zero.  The fixed divisor
$p^*N_j$ must compensate that coefficient; moreover only finitely many
such $Q$ can occur.  Taking orders yields
\[
 v(N_i)\ge v(N_j)+l_i v(D)-v(f^*R).
\]
The left side is bounded and $l_i\to\infty$, so $v(D)=0$.
The remaining task is to find base sections of small normalized order.

\subsection*{The base line is nef}

The normalized section $s_D$ on the equidimensional model has order
zero on a component over every base prime.  The preceding generic
fiber calculation gives ordinary adjoint rank one for every
$l\ge1$, with generator $s_Es_D^l$ in base frames.  Thus all the
hypotheses of Lemma~\ref{V:bounded-maximum} hold with $L=A$,
$\tau=s_{D^V}$ and the base line $B$.  It follows that $B$ has a
semipositive metric with locally bounded weights, and in particular
is nef.  The fixed divisor has zero arc order; it remains to use the
nefness of $B$ to construct base sections of small normalized order.

\subsection*{Nef reduction and a linearly descended base divisor}

The base has $q(Y)=0$: pull holomorphic one-forms through a resolution
of the dominant rational map $W\dashrightarrow Y$, and use birational
invariance of irregularity.  Apply Lemma~\ref{V:nef-descent} to $B$.
Lift $S$ further to a smooth model $S'$ mapping to its $Y'$, and
compose to obtain $g:S'\to Z$.  Multiply $a,D,B,L$ by the common
integer from that lemma.  Retain $D$ for its pullback, put
$A'=p^*A$ and $E'=K_{S'}-p^*K_W$, and obtain
\[
 aA'\sim D+g^*L,
 \qquad L\text{ Cartier and nef},\qquad v(D)=0.
\]
Equation~\eqref{V:global-sections-base}, the projection formula,
and $h_*\mathcal O_{Y'}=\mathcal O_Z$ give for every $l\ge0$
\begin{equation}\label{V:adjoint-sections-base}
 H^0(S',K_{S'}+(1+la)A')
   =s_{E'}s_D^l g^*H^0(Z,lL).
\end{equation}
If $Z$ is a point, the divisor $D$ already proves the theorem.
Assume henceforth that $b=\dim Z>0$.

Choose a very ample $H$ on $Z$.  For some integer $m_0>0$,
\begin{equation}\label{V:pseudoeffective-base-domination}
 P=m_0A'-g^*H\quad\hbox{is pseudoeffective}.
\end{equation}
Indeed a multiple of one interior $p^*N_i$ linearly dominates $g^*H$,
because the original image polarization is big on the modified base.
Writing $N_i\sim M+m_iA$ and adding the pseudoeffective negative
multiple of $p^*M$ proves \eqref{V:pseudoeffective-base-domination}.

\subsection*{Extending first jets and proving bigness}

Take a singular semipositive metric $h_P$ on $P$.  Skoda integrability
and a finite covering of the compact variety $S'$ give a small
rational $\epsilon>0$ with
$\mathcal I(h_P^\epsilon)=\mathcal O_{S'}$.
Use Lemma~\ref{V:sesha} with $c=b+1$ to choose $k$, a very general
$z\in Z$, the ample rational divisor
$J=\tfrac12(kL+\epsilon H)$, its smooth metric of curvature $\eta$,
and $\psi$ with
\[
 \eta+i\partial\bar\partial\psi\ge0,
 \qquad\psi(w)\le(b+1)\log|w|^2+O(1).
\]
Choose $z$ also so that some $x\in g^{-1}(z)$ lies outside $D\cup E'$
and $g$ is submersive at $x$.  The image of this nonempty open subset
of $S'$ contains a dense base open, so the additional condition is
compatible with being very general.

For all sufficiently large $l\ge k$ put $m=1+la$ and use
\begin{equation}\label{V:jet-metric-decomposition}
 mA'\sim_{\mathbb Q}
 kD+g^*(2J)+\epsilon P+(m-ka-\epsilon m_0)A'.
\end{equation}
The last coefficient is nonnegative.  The divisor metric of $kD$,
the smooth metric of $2J$, the metric $h_P^\epsilon$, and the smooth
semipositive metric on the remaining multiple of $A'$ define a
singular metric $h_m$ on the integral line bundle $mA'$.  Rational
identifications are interpreted on a common multiple and then rooted.
Vanishing by $kD$ cancels its divisorial singularity, so
\[
 \mathcal O_{S'}(-kD)\subseteq\mathcal I(h_m).
\]
Put $\Psi=g^*\psi$.  Locally $\psi=c\log\sum_j|f_j|^2+u$
with $c>0$ and $u$ smooth; pullback has the same form with
$f_j\circ g$ and $u\circ g$.  Surjectivity of $g$ prevents it
from being identically $-\infty$, so it still has neat analytic
singularities.  The two copies of the base curvature in
\eqref{V:jet-metric-decomposition} give
\[
 i\Theta_{h_m}+(1+s)i\partial\bar\partial\Psi\ge0
 \qquad(0\le s\le1).
\]
The extension theorem of Cao--Demailly--Matsumura
\cite[Theorem~1.1]{CDM2017} applies on the compact K\"ahler manifold
$S'$: $h_m$ may be singular, and $\Psi$ has neat analytic
singularities.  Injectivity in degree one, together with the long
exact sequence, yields the surjection
\begin{equation}\label{V:quotient-surjection}
 H^0\!\left(S',\mathcal O_{S'}(K_{S'}+mA')\otimes\mathcal I(h_m)\right)
 \longrightarrow
 H^0\!\left(S',\mathcal O_{S'}(K_{S'}+mA')\otimes
 \frac{\mathcal I(h_m)}{\mathcal I(h_me^{-\Psi})}\right).
\end{equation}

The quotient is supported over $z$, since $\Psi$ is smooth elsewhere.
For any holomorphic local section $\sigma$ of $lL$ near $z$, the
section $s_{E'}s_D^l g^*\sigma$ lies in the larger ideal twist near
the entire fiber, because $l\ge k$.  Its quotient class glues with
zero off the fiber.  Surjectivity and
\eqref{V:adjoint-sections-base} produce $u\in H^0(Z,lL)$ for which
\[
 s_{E'}s_D^l g^*(u-\sigma)
\in \mathcal O_{S'}(K_{S'}+mA')\otimes\mathcal I(h_me^{-\Psi})
\]
near $x$.  The effective factors are units at $x$, and the local weight
of $h_m$ is bounded above.  Submersion coordinates and Fubini's
theorem therefore imply local integrability downstairs of
\[
 |(u-\sigma)(w)|^2|w|^{-2(b+1)}.
\]
If the first nonzero homogeneous term had degree $d\le1$, integration
over a cone of directions where it does not vanish would contain
$\int_0^\varepsilon r^{2d-2(b+1)+2b-1}\,dr$, which diverges.
Thus $u-\sigma$ vanishes to order at least two.  First jets of $lL$
are generated at $z$.

Choose a local frame and prescribe the jets $1,w_1,\ldots,w_b$.
The resulting global sections define ratios with independent
differentials at $z$, so the associated rational map has image
dimension $b$.  Consequently $L$ is big.

\subsection*{Finishing the order estimate}

Write $L\sim_{\mathbb Q}H_1+C$ with $H_1$ ample rational and $C\ge0$.
For every sufficiently small rational $\gamma>0$, the rational divisor
$(1-\gamma)L+\gamma H_1$ is ample.  A sufficiently divisible basepoint
free multiple has a section not vanishing at the center of the lifted
arc on $Z$, even if this center is not a closed point.  Indeed global
generation gives a section which is a unit in the local ring at that
scheme point.  Dividing its divisor by the multiple gives an effective
rational divisor $H_\gamma\sim_{\mathbb Q}(1-\gamma)L+\gamma H_1$
with $v(g^*H_\gamma)=0$.  Hence
\[
 D+g^*(H_\gamma+\gamma C)\sim_{\mathbb Q}aA',
 \qquad v\bigl(D+g^*(H_\gamma+\gamma C)\bigr)
       =\gamma v(g^*C).
\]
After clearing denominators this is the divisor of a section of a
positive multiple $rA'$; such sections descend to $W$.
Its order divided by $r$ is $\gamma v(g^*C)/a$, which is less than
$\delta$ for a suitable $\gamma$.  This completes the proof.
\end{proof}
